% Generated draft: requires mathematical and source review.
\documentclass[11pt]{article}
\usepackage[margin=25mm]{geometry}
\usepackage{fontspec}
\setmainfont{Latin Modern Roman}
\usepackage{amsmath,amssymb,mathtools,unicode-math}
\setmathfont{Latin Modern Math}
\usepackage{enumitem,needspace,xurl,hyperref}
\setlength{\parindent}{0pt}
\setlength{\parskip}{5pt}
\setlength{\emergencystretch}{3em}
\newcommand{\N}{\mathbb{N}}
\newcommand{\Z}{\mathbb{Z}}
\newcommand{\Q}{\mathbb{Q}}
\newcommand{\R}{\mathbb{R}}
\newcommand{\C}{\mathbb{C}}
\newcommand{\F}{\mathbb{F}}
\newcommand{\A}{\mathbb{A}}
\newcommand{\PP}{\mathbb{P}}
\newcommand{\E}{\mathbb{E}}
\newcommand{\eps}{\varepsilon}
\newcommand{\dd}{\,\mathrm{d}}
\begin{document}
\section*{UnsolvedMath ID: 20000864}
\textbf{Source problem code:} catalogue entry
\subsection{Definitions and mathematical statement}
\textbf{Nature of Chapter G.} This is a discussion after Ken Nakayama's talk, not a formal conjecture. Dana asks whether generic views agree with Bayes. The reply calls Bayesian priors stochastic/nonergodic, invokes Kanizsa diagrams and criticizes Bayesian unconstrainedness. Other comments ask whether generic views themselves are priors, mention Bill Freeman's Bayesian treatment, note ambiguous figures (speaker marked ``Jeff?''), and ask about motor experience. The discussion contrasts impossible-triangle perception with explicit knowledge, asks about neurobiological explanation and models versus theory, records ``Reductionism is metaphysical,'' and emphasizes developmental kinetic-depth/stereo trajectories and perception--motor interaction. Its optimistic/pessimistic-camp remark is commentary, not a mathematical assertion.

The aid's identified mathematical issue is whether generic-view accounts are distinguishable from unconstrained Bayesian rationalizations. The following finite-model statements are introduced by that aid; they are not equations or posed theorems in the transcript.

\textbf{Reported prior-absorption criterion.} Let $\mathcal H,\mathcal I$ be finite hypothesis/image sets, with positive prior $\pi(h)$, positive specified likelihood $L(i\mid h)$ and positive genericity multiplier $G(h,i)$. The posterior proportional to $\pi(h)L(i\mid h)G(h,i)$ can use the same likelihood and absorb $G$ entirely into one image-independent object prior exactly when
\[
 G(h,i)=a(h)b(i)
\]
for positive functions $a,b$. Then the new prior is proportional to $\pi(h)a(h)$. Full positivity makes this equivalent to all cross-product identities $G(h,i)G(h',i')=G(h,i')G(h',i)$. With zeros, the criterion must be interpreted on the actual likelihood support; the supplied summary does not state those cases.

\textbf{Reported arbitrary-table rationalization.} Given any conditional table $Q(h\mid i)\geq0$ with $\sum_hQ(h\mid i)=1$, choose any strictly positive image evidence $e(i)$ summing to one. Define
\[
 P(h,i)=e(i)Q(h\mid i),\qquad
 \pi(h)=\sum_i e(i)Q(h\mid i),\qquad
 L(i\mid h)=\frac{e(i)Q(h\mid i)}{\pi(h)}\quad\text{if }\pi(h)>0.
\]
If $\pi(h)=0$, every $Q(h\mid i)=0$ and $L(\cdot\mid h)$ may be any probability vector. These choices give a normalized prior/likelihood model reproducing $Q$ on every image, since $e(i)>0$. The aid states that every joint law with full image support is recovered by choosing its own image marginal as $e$. For a general joint law, conditional probabilities are fixed only where the image marginal is positive; zero-evidence images remain version-arbitrary. Zero-prior hypotheses and zero-evidence images are therefore different phenomena.

\textbf{Reported common-prior intervention test.} Suppose environments $v$ implement specified viewpoint interventions and supply fixed positive generic-view kernels $K_v(i,h)$, incorporating any declared nuisance-view integration and genericity factors. Given positive observed posteriors $Q_v(h\mid i)$, the existence of a common positive prior with
\[
 Q_v(h\mid i)=\frac{\pi(h)K_v(i,h)}{\sum_{h'}\pi(h')K_v(i,h')}
\]
is equivalent to constancy across images and environments of every corrected odds ratio
\[
 \frac{Q_v(h\mid i)}{Q_v(h'\mid i)}\frac{K_v(i,h')}{K_v(i,h)}
 =\frac{\pi(h)}{\pi(h')}.
\]
The kernels and interventions must be specified independently; arbitrary refitting removes this restriction. The summary reports a worked two-view surface family with a prior-free ranking-reversal prediction, but gives no numerical kernels or reversal values, so that example is not reconstructed here.

The aid attributes to Nakayama--Shimojo a generic-sampling account as a degenerate Bayesian form neglecting object priors, and to Freeman viewpoint-prior integration yielding a marginal image likelihood. Full primary references and the transcript's nonergodicity definitions are absent. These conditional finite-model audits do not settle the discussion's philosophical, developmental or neural-mechanism questions. The registry says open/\texttt{OPEN-TRIAGE}, while the attempt itself says context-only/not-a-problem; both scopes are retained.
\subsection{Short English statement}
The source is a 2003 discussion about generic views, Bayes and neural explanations. The aid supplies finite-model tests: when a genericity factor can become an object prior, why an unrestricted conditional table always has a Bayesian rationalization, and how fixed viewpoint interventions can test a common-prior restriction. These formalizations do not turn the transcript into a solved formal conjecture.
\subsection{Sources}
\begin{itemize}
\item Local UnsolvedMath registry, permanent ID $20000864$, \texttt{catalogue entry}; canonical \texttt{aim-biology-notes.json}, \texttt{notes[28]}, tagged section.
\item AIM, ``Inference and Prediction in Neocortical Circuits,'' Chapter G, ``Helmholtz Inference in Early Vision Areas,'' discussion after Ken Nakayama's talk: \url{https://aimath.org/WWN/brain/brain.pdf}. Supplied verification dates the hard-copy web-page version October 24, 2003 and locates it at PDF indices 7--8, printed pages 8--9; it contains no talk-summary equations or formal conjecture.
\item Transcript names Dana, Jeff (one attribution uncertain), and Bill Freeman; the aid also names Nakayama--Shimojo and Freeman's primary lineage. No full paper titles or bibliographic entries are supplied, and no independent source visit is claimed.
\item Research attempt 1, \texttt{context\_only}, run status \texttt{not\_a\_problem}, internal rigor \texttt{proved}, low novelty confidence. Explicitly a machine-generated research aid, not peer reviewed; independently verify before citation. Original literature field empty.
\end{itemize}
\end{document}
